\documentclass[11pt,a4paper]{article}
\usepackage[margin=25mm]{geometry}
\usepackage{amsmath,amssymb}
\usepackage[hidelinks]{hyperref}
\input{glyphtounicode}
\pdfgentounicode=1
\setlength{\parindent}{0pt}
\setlength{\parskip}{7pt}
\title{A conditional stationary-headway interpretation of H1's capacity curve}
\author{VSL task-definition review: Step 3 supporting note}
\date{14 October 2026}
\begin{document}
\maketitle

\textbf{Status.} Source-backed stationary calculation and original algebraic
correspondence. No vehicle trajectory, traffic experiment, parameter fit or
new H1/T1 law is supplied. The approved model is unchanged. This note supports
the plain-language reference-qualification report, not an empirical capacity
or gradient claim.

\section*{1. Why speed can change usable lane density}

The official SUMO Road Capacity documentation gives the minimum gross
front-to-front spacing for homogeneous Krauss vehicles travelling at the
same constant speed. Let vehicle length be $l_{\rm veh}$ in metres,
standstill gap $g_{\min}$ in metres and desired following time $\tau_d>0$
in seconds. Write $d=l_{\rm veh}+g_{\min}>0$. At actual speed $s>0$,
the documented stationary relation and its ideal flow consequence are
\begin{align}
 D(s)&=d+\tau_d s, & [D]&={\rm m},\\
 k_{\rm ideal}(s)&=\frac{1}{d+\tau_d s},
 & [k_{\rm ideal}]&={\rm veh/m},\\
 Q_{\rm ideal}(s)&=\frac{s}{d+\tau_d s},
 & [Q_{\rm ideal}]&={\rm veh/s}.
\end{align}
The documentation cautions that these conditions are rarely achieved
exactly in a simulation. Desired-speed differences, random slowdowns and
nonminimum gaps are among its stated reasons. This is an ideal stationary
relation, not the realised capacity of the proposed lane drop.

\section*{2. Algebraic correspondence with the accepted H1 curve}

H1 uses triangular-diagram critical density and capacity
\begin{equation}
 k_c(v)=\frac{w k_j}{v+w},\qquad C(v)=v k_c(v).
\end{equation}
For the sole purpose of interpreting the stationary curve, set
\begin{equation}
 k_j=\frac{1}{d},\qquad w=\frac{d}{\tau_d}.
\end{equation}
Substitution gives
\begin{equation}
 k_c(v)=\frac{1}{d+\tau_d v},\qquad
 C(v)=\frac{v}{d+\tau_d v}.
\end{equation}
Thus the two curve shapes coincide if the stationary actual speed $s$
equals the predictor's advice $v$. This equality of shapes is an inference
from the documented headway law. It does not establish equality of the two
dynamic traffic models or prescribe reference parameters.

\newpage
\section*{3. The traffic meaning and its limits}

The directions of the stationary changes are
\begin{equation}
 \frac{d k_{\rm ideal}}{d s}
 =-\frac{\tau_d}{(d+\tau_d s)^2}<0,
 \qquad
 \frac{d Q_{\rm ideal}}{d s}
 =\frac{d}{(d+\tau_d s)^2}>0.
\end{equation}
Slower travel permits more closely spaced vehicles per metre at the same
desired following time, while ideal vehicles-per-second throughput falls.
This gives a physical interpretation to density headroom in H1 without
claiming that lower speed increases every kind of service. In H1's progress
coordinate, the explicit speed factor in the preparation law cancels:
the remaining preparation effect comes from headroom and the model's
specified interaction, not additional seconds by themselves.

The following conclusions do \emph{not} follow from the correspondence:
\begin{enumerate}
 \item A native lane-drop interaction obeys H1's receiving allocator,
 preparation opportunity scale or forced-service work law.
 \item $w=d/\tau_d$ is a measured transient backward-wave speed in the
 chosen microscopic model. The equality identifies a stationary curve only.
 \item Native actual speed equals the permitted advice during acceleration,
 braking, lane changes or queueing.
 \item A positive analytical work-factor range is traffic-calibrated, or
 an interior optimal advice exists.
\end{enumerate}
For example, the earlier analytical family chose $w=2u$. Under this optional
stationary interpretation it would imply $u=d/(2\tau_d)$. Its dimensionless
construction must therefore not be presented as already calibrated freeway
speeds. The general approved family is broader; no numerical mapping is
selected here.

\section*{4. Why a mixed departure trace may not identify the forced server}

On intervals where the accepted fluid server is binding, let $u_A$ be
continuing departure flow and $g_E$ ending departure flow, both in
vehicle-equivalents per second. Its work account has the diagnostic form
\begin{equation}
 u_A+\alpha g_E=R.
\end{equation}
One repeated mixed-flow value supplies only one relation for the two
unknowns $\alpha$ and $R$. At least differing flow observations with a
nondegenerate design, correct binding-interval classification and a valid
common-law fit are needed to separate them; two different $g_E$ values are
necessary for rank two in this simple linear diagnostic. Native stopped
vehicles do not by themselves certify the fluid server's binding condition.
If the law fails across independently prepared development traces, or
parameters remain unidentifiable, retain that failure or the admissible
bounds. Do not select $\alpha$ from the positive-action range or fit held-out
optima. This is an identifiability observation, not a fitted native model.

\textbf{Primary source inspected, 14 October 2026:}
\href{https://sumo.dlr.de/docs/Simulation/RoadCapacity.html}{SUMO documentation,
Road Capacity}, sections Computing Headways and Further headway effects.
All correspondence, derivatives and identifiability reasoning above are
the review's deductions under their stated conditions.
\end{document}
